\documentclass{amsart}

% 数学相关宏包
\usepackage{amsmath,mathtools,amsthm,amssymb}
\usepackage{bbm}
\usepackage{mathrsfs}
\usepackage{centernot}

% 图形和绘图
\usepackage{graphicx,float}
\usepackage{tikz,tikz-cd}
\usepackage{pgfplots}
\usepackage[framemethod=TikZ]{mdframed}
\usetikzlibrary{decorations.pathmorphing,positioning,arrows,automata,calc,shapes,patterns}
\pgfplotsset{compat=1.18}
\usepgfplotslibrary{statistics}

% 表格和列表
\usepackage{booktabs,multirow,colortbl}
\usepackage{enumitem}
\setlist[itemize,enumerate,description]{noitemsep}
\setlistdepth{9}

% 算法
\usepackage{algorithm,algorithmic}

% 字体和颜色
\usepackage{xcolor,listings}
\usepackage{xeCJK}
\setCJKmainfont{Noto Serif CJK SC}

% 页面和链接
\usepackage{fancyhdr,caption}
\usepackage{hyperref}
\usepackage{cleveref}

% 工具宏包
\usepackage{etoolbox,letltxmacro,docmute}
\usepackage{gensymb}

% 自定义设置
\renewcommand{\arraystretch}{1.5}
\let\originalmathbb\mathbb
\renewcommand{\mathbb}[1]{%
    \ifstrequal{#1}{1}{\mathbbm{#1}}{\originalmathbb{#1}}%
}

% 颜色定义
\definecolor{lightblue}{RGB}{173,216,230}
\definecolor{lightgreen}{RGB}{144,238,144}
\definecolor{lightyellow}{RGB}{255,255,224}
\definecolor{lightgray}{RGB}{211,211,211}

		% 超链接设置
		\hypersetup{
		    colorlinks=true,
		    linkcolor=red,
		    filecolor=magenta,
		    urlcolor=cyan,
		    pdftitle={\@title},
		    pdfauthor={\@author},
		    pdfsubject={数学笔记},
		    pdfkeywords={数学},
		    bookmarksnumbered=true,
		    bookmarksopen=true,
		    bookmarksopenlevel=6,
		    pdfstartview=FitH,
		    pdfpagemode=UseOutlines,
		    pdfborder={0 0 0}
		}

% 定理环境
\theoremstyle{plain}
\newtheorem{theorem}{Theorem}[section]
\newtheorem{lemma}[theorem]{Lemma}
\newtheorem{corollary}[theorem]{Corollary}
\newtheorem{proposition}[theorem]{Proposition}

\theoremstyle{definition}
\newtheorem{definition}[theorem]{Definition}
\newtheorem{example}[theorem]{Example}
\newtheorem{exercise}[theorem]{Exercise}
\newtheorem{problem}[theorem]{Problem}

\theoremstyle{remark}
\newtheorem{remark}[theorem]{Remark}
\newtheorem{note}[theorem]{Note}
\newtheorem{observation}[theorem]{Observation}

% 解答环境
\newtheorem*{solution}{Solution}
\newtheorem*{proof*}{Proof}

% 图片路径
\graphicspath{{F:/资料/2025-summer/figures/}}

\author{乐绎华, Yue Yihua}
\address{中山大学, Sun Yat-sen University}
\email{yueyh@mail2.sysu.edu.cn}
\date{\today}
\title{Mathematical Notes}

\begin{document}

\maketitle
\tableofcontents

\section{Introduction of Preliminaries}

Let $\varphi:\Omega \subseteq \mathbb{R}^{n}\to \mathbb{R}^{m}$, $\boldsymbol{\varphi}(\boldsymbol{x})=(\varphi^{1}(x),\dots,\varphi^{m}(x))$. The Jacobian derivative is
\[
D\boldsymbol{\varphi}(x)=\begin{bmatrix}
\partial_1\boldsymbol{\varphi}(x) & \partial _{2}\boldsymbol{\varphi}(x) & \dots & \partial _n\boldsymbol{\varphi}(x)
\end{bmatrix}\in M^{m\times n}
\]
Rank-one matrices: for $a\in \mathbb{R}^{m}$, $n\in \mathbb{R}^{n}$,
\[
e_1\otimes e_1=\begin{pmatrix}
1 & 0 & \dots & 0 \\
0 & 0 & \dots & 0 \\
\vdots & \vdots &  & \vdots \\
0 & 0 & \dots & 0
\end{pmatrix}
\]
\[
a\otimes n=(a_i\cdot n_j)_{i,j}^{m,n}
\]
\begin{exercise}
Let $v\in \mathbb{R}^{n}$, $(a\otimes n) v= \left< n,v \right>a$.
\end{exercise}
\begin{definition}[differential inclussion]
Given $K\subseteq M^{m\times n}, A\in M^{m\times n}$
	\begin{enumerate}
		\item $\left.\varphi\right|_{\partial \Omega}=\mathbf{A}\mathbf{x}+\mathbf{b}$
		\item $D\varphi (x)\in K$, a.e. $x\in \Omega$.
		\item $+$ Regularily on $\varphi$.
	\end{enumerate}
\end{definition}
\subsection{Regularity?}

\begin{itemize}
	\item Control on the oscillation
	\begin{itemize}
		\item Holder: $\lvert \varphi(x)-\varphi(y) \rvert\leq C\lvert x-y \rvert ^{\alpha}$, where $0<\alpha<1$.
		\item Lipschitz:  $\lvert \varphi(x)-\varphi(y) \rvert\leq C\lvert x-y \rvert$.
	\end{itemize}
	\item Bounds on the derivative
	\begin{itemize}
		\item $W^{k,p}(\Omega,\mathbb{R}^{m})$, the sobolev spaces.
		\item Approx $\int_{\Omega}^{} \lvert D^{\gamma} \varphi\rvert ^{p} \, \mathrm{d}x<\infty$, for $0\leq\gamma\leq k$.
		\item Warning: what means derivative?
		\begin{itemize}
			\item Derivatives are understood in the weak sense.
		\end{itemize}
	\end{itemize}
\end{itemize}

Textbooks:

\begin{itemize}
	\item Evans
	\item Eienmen?
	\item Adams
\end{itemize}

\begin{definition}[piecewise affine]
$\boldsymbol{\varphi}:\Omega\to \mathbb{R}^{m}$ is called pircewise affine (pa) if $\exists \{ \Omega _i \}^{\infty}_{i=1}$, s.t.
\[
\boldsymbol{\varphi}(\mathbf{x})=\sum(\mathbf{A}_i \mathbf{x}+\mathbf{b}_i)\chi_{\Omega _i}
\]
Where $m\left( \Omega\setminus \bigcup_{i=1}^{\infty}\Omega _i \right)=0$, $D\boldsymbol{\varphi}=\sum \mathbf{A}_i\chi_{\Omega _i}$.
\end{definition}
Then
\[
\int \lvert D\varphi \rvert ^{p}=\sum \lvert \Omega _i \rvert \cdot \lvert \mathbf{A}_i \rvert ^{p}
\]
We know that $W^{1,\infty}=\text{Lip}$, and $W^{1,p}\overset{ \text{inclusion} }{ \to } C^{1-\frac{n}{p}}$ for $p>n$, which is Holder space.

\subsection{History}

\begin{itemize}
	\item Nash: Isometric inmersions
	\item Approximate solutions
\end{itemize}
\[
I[u]\coloneqq \int_{}^{}W[Du(x)]  \, \mathrm{d}x 
\]
Where $W:M^{m\times n}\to \mathbb{R}_{+}$, $K\coloneqq W^{^{-1}}[0]$. (Mullen)

\begin{itemize}
	\item Regularity of elliplic systems
\end{itemize}


\section{IPM}

$\rho\in \mathbb{R}, u\in \mathbb{R}^{3}$, $\rho=\rho(x,t)$, $u=(u^{1}(x,t),u^{2}(x,t),u^{3}(x,t))$. where $x\in \mathbb{T}$, $t\in[0,T]$
\[
\begin{cases}
\partial _{t}u+\mathrm{div}(\rho u)=0 \\
\mathrm{div}(u)=0 \\
u=\nabla_{\text{p}}\cdot \rho(0,1)
\end{cases}
\]
A PDE: for $U\in K$, solve $\nabla \times U= \mathbf{0}$.

Use new variable $m=(m^{1}(x,t),m^{2}(x,t),m^{3}(x,t))$, then
\[
\begin{cases}
\partial _{t}\rho +\mathrm{div}(m)=0 \\
\mathrm{div}(u)=0 \\
\nabla \times(u+\rho(0,1))=0
\end{cases}
\]
This is a \textbf{linear equation} in $K^{\text{IPM}}\coloneqq \{ (\rho,m,u):m=\rho u \}\subset \mathbb{R}^{7}$. (Add three dimension to linearize it.)

\begin{definition}[gradient distribution]
$\varphi:\Omega\to \mathbb{R}^{m}$, $\left.\varphi\right|_{\partial \Omega}=\mathbf{A}\mathbf{x}+\mathbf{b}$, $V_{\varphi}\in \mathcal{M}(M^{m\times n})$ is gradient distribution.
\end{definition}
For open subset $\mathcal{O}\subset M^{m\times n}$, we define
\[
V_{\varphi}(\mathcal{O})=\frac{m(\{ x\in \Omega:D\varphi(x)\in \mathcal{O} \})}{\lvert \Omega \rvert }
\]
Then $C^{\infty}_{0}(M^{m\times n})=\mathcal{M}(M^{m\times n})$,
\[
\int_{M^{m\times n}}^{} f(x) \, \mathrm{d}V_{\varphi}=\int_{\Omega}^{} f(D\varphi(x)) \, \mathrm{d}x 
\]
Recall that for the measure $\delta_{A}$
\[
\int_{\Omega}^{} f \, \mathrm{d}\delta_{A} =f(A)
\]
Then for $V=\sum l_i\delta_{A_i}$, we have
\[
\int_{\Omega}^{} f \, \mathrm{d}V=\sum l_if(A_i)
\]
For $\varphi(x)=Ax$, we have $D\varphi=A$, then (why?) $V_{\varphi}=\delta_{A}$, then if $D\varphi=\sum A_i\chi_{\Omega _i}$, then $V_{\varphi}=\sum\frac{\lvert \Omega _i \rvert}{\lvert \Omega \rvert}\delta_{A_i}$.


\end{document}